% ============================================================
% Chapter4.tex — Implementation
% ============================================================
\chapter{Implementation}

\section{Technology Stack}

The Iris platform is implemented entirely in Python, leveraging the following key libraries and services:

\begin{table}[H]
\centering
\caption{Technology Stack}
\begin{tabularx}{\textwidth}{l l X}
\toprule
\textbf{Component} & \textbf{Technology} & \textbf{Purpose} \\
\midrule
LLM Inference   & Groq Cloud + LLaMA 3.3 70B Versatile & Intent extraction and response synthesis \\
Agent Framework & LangChain (\texttt{@tool} decorators) & Tool binding and JSON schema generation \\
UI Dashboard    & Streamlit + Custom CSS & Real-time observability interface \\
Gmail Access    & Google Gmail API + OAuth 2.0 & Email reading, drafting, and sending \\
Audit Storage   & JSON + CSV + SQLite (SQLite3) & Three-format concurrent audit persistence \\
Input Validation& Python \texttt{re} module (Regex) & Domain typo detection \\
Scheduling      & Python \texttt{threading.Thread} & Background daemon task execution \\
\bottomrule
\end{tabularx}
\end{table}

\section{Controlled Intent Router Implementation}

The orchestrator in \texttt{agent/controller.py} implements the \gls{cir} pattern through the following sequence:

\begin{enumerate}[leftmargin=1.5em]
  \item \textbf{Tool binding}: All registered tools are wrapped with LangChain's \texttt{@tool} decorator, which generates JSON schemas describing each function's arguments. These schemas are bound to the LLM via \texttt{llm.bind\_tools(tools)}.
  \item \textbf{Single-turn invocation}: The LLM is invoked once with the user query and tool schemas. It either returns a conversational message or a structured \texttt{tool\_call} object.
  \item \textbf{Policy enforcement}: If a tool call is detected, the tool name is passed to \texttt{risk\_classifier.py}, which performs an O(1) dictionary lookup to determine the risk tier.
  \item \textbf{Decision branching}: The decision engine applies the approval matrix and routes to execution, hold, or block.
  \item \textbf{Audit write}: Regardless of the outcome, \texttt{audit\_logger.py} writes a complete structured record.
\end{enumerate}

\begin{lstlisting}[caption={Core dispatch logic in agent/controller.py}, label={lst:dispatch}]
# Single-turn LLM invocation
ai_msg = _llm_with_tools.invoke(messages)

# If tool call detected
if ai_msg.tool_calls:
    tool_name = ai_msg.tool_calls[0]["name"]
    
    # Policy evaluation
    explanation = generate_explanation(tool_name)
    risk_level  = explanation["risk_level"]
    decision    = explanation["decision"]
    
    if decision == "reject":
        # CRITICAL: block immediately
        log_action(..., status="SKIPPED", decision="reject")
    elif decision == "need approval":
        # HIGH: return approval card to UI
        log_action(..., status="SKIPPED", decision="need approval")
    else:
        # LOW/MEDIUM: auto-execute
        result = execute_tool(tool_name, args)
        log_action(..., status="SUCCESS", decision="approved")
\end{lstlisting}

\section{Risk Classifier Implementation}

The risk classifier (\texttt{agent/risk\_classifier.py}) uses a static dictionary mapping every registered tool name to its security tier. The classification function performs an O(1) lookup:

\begin{lstlisting}[caption={Risk classification core logic}, label={lst:risk}]
ACTION_RISK_MAP = {
    "read_recent_emails":       "LOW",
    "send_email":               "HIGH",
    "execute_terminal_command": "CRITICAL",
    "query_database":           "HIGH",
    "buy_stock":                "CRITICAL",
    # ... 65+ more entries
}

def classify_risk(action: str, default: str = "MEDIUM") -> str:
    return ACTION_RISK_MAP.get(action.lower(), default)
\end{lstlisting}

\section{Audit Logger Implementation}

The logger (\texttt{agent/audit\_logger.py}) implements a thread-safe write-through cache strategy. An in-memory Python list (\texttt{\_log\_buffer}) acts as the primary read store for the UI, while writes are propagated to all three disk formats inside a \texttt{threading.Lock()} guard:

\begin{lstlisting}[caption={Thread-safe multi-format audit persistence}, label={lst:logger}]
_log_buffer: list[dict] = []
_buffer_lock = threading.Lock()

def log_action(...) -> dict:
    log_entry = {
        "event_id": f"EV-{uuid.uuid4().hex[:6].upper()}",
        "timestamp": datetime.now(timezone.utc).isoformat(),
        "user_instruction": user_instruction,
        "risk_level": risk_level,
        "decision": decision,
        # ... other fields
    }
    with _buffer_lock:
        _log_buffer.append(log_entry)
        _persist_logs(_log_buffer)  # JSON + CSV + SQLite
    return log_entry
\end{lstlisting}

\section{Gmail OAuth 2.0 Integration}

Email operations use the Gmail REST API with OAuth 2.0 authorization. Scoped credentials are stored in \texttt{token.json} to authorize requests without requiring password storage. The implementation supports:

\begin{itemize}[leftmargin=1.5em]
  \item Reading recent emails (subject, sender, snippet).
  \item Reading full email body by message ID.
  \item Sending emails via SMTP-over-Gmail.
  \item Creating draft emails.
  \item Applying labels to messages.
\end{itemize}

\section{Pre-Execution Typo Validator}

Before routing to the LLM, the query validator (\texttt{agent/query\_validator.py}) runs a regex-based scan to detect common spelling errors in email domains and command keywords:

\begin{lstlisting}[caption={Typo detection regex patterns}, label={lst:typo}]
EMAIL_TYPO_MAP = {
    r"\b([a-zA-Z0-9._%+-]+)@gmal\.com\b":  "gmail.com",
    r"\b([a-zA-Z0-9._%+-]+)@gamil\.com\b": "gmail.com",
    r"\b([a-zA-Z0-9._%+-]+)@gmial\.com\b": "gmail.com",
    # ... more patterns
}
\end{lstlisting}

If a typo is detected, a warning message is displayed in the dashboard \textit{before} the query proceeds to the LLM, allowing the user to correct it without executing a flawed action.

\section{Streamlit Dashboard Architecture}

The dashboard (\texttt{app.py}) is organized into three tabs:

\begin{itemize}[leftmargin=1.5em]
  \item \textbf{Chat Tab}: Main conversational interface with per-message execution trace cards showing tool, risk badge, decision badge, and policy reasoning.
  \item \textbf{Audit Log Tab}: Filterable, searchable database view with session metrics (total actions, success rate, risk distribution chart) and JSON export button.
  \item \textbf{System Tab}: Live tool registry showing all 14+ registered tools with their risk tiers and descriptions.
\end{itemize}

A custom \texttt{style.css} file (21KB) provides the full visual identity: dark-themed chat bubbles, HSL-tailored risk badges, glassmorphism metric cards, and smooth CSS transitions.
